\chapter{Conclusion} \label{chap:conclusion}

This thesis has presented a novel approach to whispered-to-voiced speech conversion using synthetic voicing techniques. The conclusions drawn from this research can be attributed to the performance of the proposed platform and the perceptual data that has been collected. The following sections summarize the contributions, limitations, and future work, along with the final remarks and acknowledgements.

First of all, the proposed system has demonstrated the feasibility of the conversion in a realtime scenario, with the ability to generate voiced speech from whispered input on-the-fly without any significant delays. The system architecture, comprising the glottal pulse generation and vocal tract filter modelling modules, has been successfully implemented and integrated into the STM32 platform. Despite the limited computational resources of the development platform, especially regarding the pulse generation, promising results have been obtained in terms of speech quality and intelligibility.

Moreover, the whispered speech detection and segmentation algorithms have proven to be effective in identifying the voiced and unvoiced regions of the input signal, allowing for the selective application of synthetic voicing, using a lightweight and efficient criteria. Despite this, this criteria is applied on a frame-by-frame basis, which may lead to discontinuities in the generated speech, especially in the presence of noise or rapid changes in the input signal from phoneme to phoneme.

Future work should focus on improving the voiced speech generation by perhaps exploring other realtime processing platforms such as PortAudio, which could provide more computational resources and allow for more sophisticated algorithms to be implemented. The integration of a more robust whispered speech detection and segmentation algorithm, capable of identifying phoneme boundaries and transitions, could also improve the quality of the generated speech. Machine learning techniques could also be employed to enhance the system's performance by providing more accurate estimates of voicing decisions and the fundamental frequency that is to be applied to the synthetic pulses. Perhaps the existing segmentation models could be improved upon by pretraining and learning the parameters for centre of gravity mapping and the WAD algorithm. This could be achieved by using the existing labeled data, which includes markings for voiced and unvoiced regions, to train a machine learning model that can predict the optimal parameters for the segmentation algorithms. This could potentially improve the accuracy of the segmentation and reduce the discontinuities in the generated speech. 

%\vspace*{12mm}